\begin{abstract}
Injection--recovery characterizes CNN strong-lens finders, but its
interpretation depends on the simulated images presented to the detector. For
a frozen EfficientNetV2-S trained on DESI Legacy Imaging Surveys DR10 cutouts,
we measure conditional recovery and observed--injected distinguishability
separately. Matched and paired controls identify differences already present
before source insertion and measure detector responses to specified
construction changes. Dependency-aware linear probes compare four
representations---the operational detector's pooled penultimate features,
self-supervised survey features, ImageNet features, and pixel statistics---for
43 selected observed systems and 6,336 source-injected images.
The populations separate in all four representations, with out-of-fold
receiver-operating-characteristic AUCs of 0.991, 0.958, 0.979, and 0.981,
respectively. A leading pixel-statistic offset is already present between
observed systems and bare targets. Along a separately fitted operational
observed--bare axis, the mean injection movement is near zero despite
heterogeneous pairwise effects. Separately, the finder recovers 3.45 per cent
of a 49,500-scene grid with no visibility cut. Modelled source-photon sampling
lowers recovery, whereas the single-S\'{e}rsic construction is recovered more
often than its paired structured-source counterpart. On 43 fixed targets,
expectation-image recovery varies non-monotonically with source brightness and
peaks at 18.6 per cent.
Together, these analyses reveal differences between the selected populations
within the actual finder, identify a pre-insertion offset, and measure
responses to controlled injection choices. They provide a practical basis for
comparing future injection pipelines using trained finders.
\end{abstract}
